\section{Updated benchmark, baselines and experimental protocol}
\subsection{Dataset and code representation}
MovieLens 1M~\cite{movielens,harper2015} contains 1,000,209 ratings, 6,040 users
and 3,706 rated movies; its metadata catalogue has 3,883 movies. SHA-256 checks
verify the official archive and all three input files. Rating density is
4.468\%; absent user--movie pairs are unobserved, never zero-star labels.

\begin{table}[!htbp]\centering\small
\begin{tabularx}{\linewidth}{@{}p{.23\linewidth}XX@{}}\toprule
File / field & Meaning & Representation / use\\\midrule
ratings: UserID, MovieID & Raw entity identifiers & Zero-based arrays \code{u}, \code{i}; mappings \code{uid}, \code{iid}\\
ratings: Rating & Observed integer, 1--5 & Target array \code{r}; train/validation/test separated\\
ratings: Timestamp & Unix observation time & Audited; no temporal predictor or chronological claim\\
movies: MovieID, Title & Identifier and display name & ID joins; titles are not numeric features\\
movies: Genres & Pipe-separated genre labels & 18-column multi-hot matrix used only for graph construction\\
users: UserID & User identifier & Join/audit only beyond the rating ID\\
users: Gender, Age & Recorded demographic category, age bin & Parsed/audited; excluded from prediction\\
users: Occupation, Zip-code & Occupation category and postal code & Parsed/audited; excluded from prediction\\\bottomrule
\end{tabularx}
\caption{All dataset fields and their actual uses. IDs are categories, not numerical magnitudes.}
\end{table}

At $K=8$, each user and each movie has eight unnamed learned traits plus one
bias: $(6040+3706)(8+1)=87,714$ persistent scalar latent parameters.
The eight coordinates are not predefined genres. EP stores two natural
parameters per scalar plus per-rating sites. GMRF adds fixed graph weights;
the count above excludes observations, deterministic auxiliary variables,
hyperparameters, and EP site-cache entries.

\subsection{Common evaluation and validation-only selection}
Five seeds (42--46) each define 800,167 training, 100,020 validation and
100,022 test rows through one seeded random permutation. All compared methods
use the same rows within a seed. Seeds control split/order and initialization;
42 has no special statistical meaning. The first three splits reproduce the
M2 protocol. Results from this fresh execution are reported together; old M2
timings are not mixed with current timings.

\textbf{Baseline families.} Global/user/item frequency predictors use the M2
smoothed class-histogram rules. MF-SGD retains the M2 kernel, $K=8$,
learning rate 0.007, epoch decay 0.98, regularization 0.03 and 25 epochs.
MF-MAP and GMRF-MAP use the same alternating solver with and without the
graph factor. The two MF optimizers use different objectives/regularization
conventions, so their numerical regularization values are not interchangeable.

\textbf{EP family.} ADF and EP1--5 share the original priors and hybrid local
updates; EP10 studies a larger budget. EP1, EP2, etc. mean the number of passes,
not different models. An additional experiment fits EP5 at observation
variances $\beta^2\in\{0.5,1,2\}$; validation NLL selects the noise candidate.
All fixed candidates are also reported to reveal sensitivity.

\textbf{Graph family.} The clean-data search uses
$\gamma\in\{0,1,10,100\}$ and 25 epochs per candidate. Validation clipped-score
RMSE selects the checkpoint and graph strength. \code{GMRF\_MAP} reports the
best positive strength as a graph diagnostic; \code{Graph\_selected} permits
$\gamma=0$, so a deployment choice can switch the graph off. MF-MAP is the
matching $\gamma=0$ ablation. Every candidate uses the same initialization and
epoch budget. The Gaussian prediction scale is then fitted on validation NLL.
No test metric enters these selection functions.

\subsection{Robustness, ablations and noise}
Three seeds (42--44) are used for controlled stress tests, with validation and
test labels kept clean. Two different meanings of noise are examined:
\emph{assumed observation variance} inside EP and \emph{corruption of observed
training labels}. They answer different questions.

\begin{table}[!htbp]\centering\small
\begin{tabularx}{\linewidth}{@{}lXX@{}}\toprule
Experiment & Change & Fixed comparison / interpretation\\\midrule
Training availability & Retain 25\%, 50\%, 100\% of training rows & Nested seeded subsets; identical held-out rows; random missingness only\\
Label noise & Replace exactly 10\% or 20\% of training ratings by a uniformly selected different star & Same corruption for every method; graphs rebuilt from corrupted training data\\
Graph density & Top 5, 20 or 50 outgoing neighbours & Same mixture/regularization; symmetrized final degree can be larger\\
Feature ablation & Ratings only, genres only, mixed graph & Distinguishes information sources; no test-derived similarities\\
Entity cold-start & Separately remove all training rows of 10\% of item/user identities & Evaluate only those entities' original test rows; item genres remain available\\
Inference sensitivity & Prior variance 0.1/0.5, damping 0.25/1, reshuffled order & EP5 against original settings, three seeds; no claim of exhaustive tuning\\\bottomrule
\end{tabularx}
\caption{Controlled experiments. Clean curves use the same three seeds as their stress conditions.}
\end{table}

For stress tests, the graph strength is fixed at 10, the candidate retained from
the clean-validation comparisons, and compared with $\gamma=0$; checkpoints
still use validation within 25 epochs. This tests robustness of that operating
setting rather than repeatedly optimizing the full grid on every condition.
The cold-start checkpoint criterion uses the full original validation
split, while its reported test metric uses only withheld items or users.
Missingness is uniform over observed training ratings, not a model of
nonrandom exposure. No ratings are generated to fill missing entries.

\subsection{KPIs, confidence intervals and reproducibility}
For probabilities $p_{nr}$, $\hat r_n=\sum_{r=1}^5rp_{nr}$.
\begin{align}
\mathrm{RMSE}&=\sqrt{N^{-1}\sum_n(\hat r_n-r_n)^2},&
\mathrm{MAE}&=N^{-1}\sum_n|\hat r_n-r_n|,\\
\mathrm{NLL}&=-N^{-1}\sum_n\log p_{n,r_n},&
\mathrm{Brier}&=N^{-1}\sum_n\sum_{r=1}^5(p_{nr}-\mathbf1\{r=r_n\})^2.
\end{align}
Top-class accuracy compares $\arg\max_r p_{nr}$ with the observed star, not a
rounded expectation. ECE uses ten equal-width top-confidence bins and their
weighted absolute confidence--accuracy gaps. Lower RMSE, MAE, NLL, Brier and
ECE are better. RMSE 0.85 is an error in star units, not 85\% accuracy.
We also retain 90\% equal-tail interval coverage/width and prediction time.

Main tables report mean $\pm$ sample SD across five seeds; robustness uses
three. Since random splits overlap, we do not treat five seed scores as five
independent datasets or use a naive t-test to certify significance. Instead,
within each split, 2,000 paired bootstrap replicates resample users with
replacement, retaining all that user's test ratings for both models. Intervals
refer to the difference in rating-weighted RMSE (baseline minus comparison).
They are conditional on fitted models, do not include retraining uncertainty,
and do not account for dependence through shared items. Several comparisons
are exploratory; no familywise discovery claim is made.

No true movie-dependency graph is available, so SHD/SID and edge-recovery
precision are not meaningful target KPIs here. We report graph size/density,
predictive error, probability quality and energy instead. A small known
Gaussian model separately permits exact inference error comparisons.

Execution uses sequential CPU kernels, numerical threads set to one, and JIT
warm-up before timing. Software versions, configurations, split hashes, source
hashes, full logs and raw result tables are supplied. MAP/SGD fit times include
validation checkpointing and scale fitting; EP fit times exclude validation
scoring. Graph construction and complete search cost are reported separately.
Thus wall times describe this implementation and machine, not an algorithmic
speed theorem. Regenerating all runs and figures is documented in the Code README.
